\section{The commission test, run on the OPCW's team and the Gaza commission}
\label{sec:B.2}
The commission test of \ref{sec:B.1} is run here on the Organisation for the Prohibition of Chemical Weapons' Investigation and Identification Team, which attributes chemical attacks in Syria. The Conference of States Parties adopted its mandate on 27 June 2018 by 82 votes to 24, over the protest of Russia, which said the move went beyond the organization's mandate \autocite{opcw2018attribution,bbc2018opcwblame}.

\emph{The threshold.} Met. The mandate was adopted by vote over the protest of Russia, Syria's principal ally, so no party to the conflict controlled the decision; the vote record is what shows it.

\emph{Evidence and methodology published.} Met. Its five reports are public in the organization's six languages, and its methods, its rules on samples and chain of custody, and its standard of proof, “reasonable grounds,” the standard fact-finding missions and commissions of inquiry use, are published \autocite{opcw2026iitpage}.

\emph{A reply offered and recorded.} Met in substance. Syria's former government didn't cooperate with any investigation before 2025, and the organization records that it considered that government's positions and found no concrete information supporting them \autocite{opcw2026iitkafrzeita}; Syria's and Russia's statements on the first report are published beside everyone else's. Whether draft findings were sent for reply before each report is in the reports, which this translation didn't read.

\emph{Displaced by a later merits judgment.} Met in design and never exercised. The team disclaims judicial status and individual responsibility \autocite{opcw2026iitpage}, and no court has reached the merits of any incident it attributed: the French, German, and Swedish proceedings concern other attacks, and the Court's case against Syria is under the Torture Convention \autocite{osji2021frenchcw,icj2023canadasyria}.

\emph{Mandate fixed before the finding.} Met on the wording: the mandate dates from June 2018 and the first finding from April 2020 \autocite{opcw2020iitltamenah}. It is contestable in spirit, since the mandate reaches only Syria and was adopted after most of the incidents it would examine. But a rule requiring the mandate to precede the events would fail nearly every commission of inquiry, since they are created in response to events, and the mandate's generic terms have produced a finding against a non-state group as well as against the state \autocite{opcw2024iitmarea}.

On the commission test, the team counts: its attributions, Douma among them \autocite{opcw2023iitdouma}, are findings by a commission of inquiry, and a merits judgment of a court would displace them if one came. Counting them licenses nothing under the force term, which lifts only to repel an armed attack under way, for force the Security Council authorizes, to stop a genocide under way, or to end an occupation found unlawful (\S\ref{sec:6.1}). A finding that a government used chemical weapons is none of these.

The same test, run on the commission of inquiry whose finding of September 2025 \S\ref{sec:6.8} rests on \autocite{coiopt2025genocide}. \emph{The threshold}: met. The Human Rights Council mandated it on 27 May 2021 by 24 votes to 9, with 14 abstentions, and neither Israel nor Palestine had a vote \autocite{ungeneva2021coi}. \emph{Evidence and methodology published}: met. The report states its standard, “reasonable grounds to conclude,” and footnotes its sources \autocite[para.~7]{coiopt2025genocide}; the Commission had no access to Gaza. \emph{A reply offered and recorded}: met in the reports the finding builds on. One of them records four requests to Israel for information and six for access, none answered \autocite[para.~2]{coiopt2024crp4}. Whether a draft of the genocide analysis was sent for reply is in pages this translation didn't read. \emph{Displaced by a later merits judgment}: met. The report says the Court “will make a final determination in the future” \autocite[para.~6]{coiopt2025genocide}. \emph{Mandate fixed before the finding}: met, and before the events as well. Israel and the United States dispute the body and not the dates. Israel's Foreign Ministry called its members “Hamas proxies” and said the report “relies entirely on Hamas falsehoods” \autocite{karmel2025unpanel}; the State Department has objected to the “open-ended and vaguely defined nature” of the mandate \autocite{nahmias2022coi}. Those objections go on the record beside the finding, as \ref{sec:B.1} requires, and the term's answer to them is the two-body rule: no commission's finding lifts it alone.

The fact-finding mission of 2009, on which \S\ref{sec:6.8} dates the occupation exception for Gaza, meets the test on the wording, with one dispute to record: the Council's resolution named Israel's violations only, the mandate covering all sides came from the Council President's letter, and Israel answered that no individual's statement could change the resolution \autocite[Annex~II]{hrc2009gaza}.
